

\documentclass[a4paper,12pt]{letter}

\usepackage[utf8]{inputenc}
\usepackage[english,slovene]{babel}

\usepackage[tikz]{bclogo}
\usepackage{qrcode,marvosym}
\usepackage{pgf,tikz,pgfplots}
\pgfplotsset{compat=1.14}
\usepackage{mathrsfs}
\usetikzlibrary{arrows}
\usepackage{graphicx}
\usepackage{fancybox}
\usepackage{amsfonts}
\usepackage{amsmath}
\usepackage{amssymb}
\usepackage{wallpaper}
\usepackage{hyperref}
\usepackage{array}
\usepackage{gensymb}
\usepackage{mdframed}
\usepackage{multicol}
\usepackage{eurosym}
\usepackage{tikz-3dplot}
\usepackage{titlesec}
\usepackage{venndiagram}
\usepackage[table]{xcolor}
\usepackage[printwatermark]{xwatermark}
\usepackage{fancyhdr}
\usepackage{enumerate}

\newcommand{\ds}{\displaystyle}
\newcommand{\ora}{\overrightarrow}

\newcolumntype{M}[1]{>{\centering\arraybackslash}m{#1}}
\newcolumntype{N}{@{}m{0pt}@{}}

\textwidth 20 true cm
\oddsidemargin -2 true cm
\evensidemargin -2 true cm
\topmargin -3.5 true cm
\textheight 27.5 true cm
\linespread{1.2}

\definecolor{bgblue}{RGB}{0,0,253}
\definecolor{zzuxux}{rgb}{0,0,0}

\newcounter{tocke}
\setcounter{tocke}{0}

\mdfdefinestyle{theoremstyle}{
linecolor=brown!40,
linewidth=1pt,
frametitlerule=true,
frametitlebackgroundcolor=brown!50,
innertopmargin=\topskip,
}

\mdfdefinestyle{theoremstyle1}{
linecolor=brown,
middlelinewidth=1pt,
frametitlerule=true,
apptotikzsetting={
\tikzset{
mdfframetitlebackground/.append style={
shade,
left color=black!40,
right color=gray!10},
mdfbackground/.append style={
top color=white,
bottom color=black!5}
}},
frametitlerulecolor=gray,
frametitlerulewidth=2pt,
innertopmargin=\topskip,
innerbottommargin=\topskip,
}

\mdtheorem[style=theoremstyle1]{naloga}{Naloga}
\mdtheorem[style=theoremstyle]{resitev}{Analiza naloge}
\mdtheorem[style=theoremstyle]{kriterij}{Kriterij ocenjevanja}

\newwatermark[
allpages,
color=brown!2,
angle=45,
scale=5,
xpos=-5,
ypos=0]{matej.info}

\titleformat{\subsection}
{\normalfont\Large\bfseries\color{brown}}
{\thesubsection}{1em}{}

\titleformat{\subsubsection}
{\normalfont\large\bfseries\color{brown}}
{\thesubsubsection}{1em}{}

\begin{document}

\pagestyle{fancy}
\renewcommand{\headrulewidth}{0pt}
\lhead{}
\chead{}
\rhead{}
\rfoot{\vspace*{-2\baselineskip}}

\renewcommand{\vec}{\overrightarrow}
\everymath{\displaystyle}

\begin{minipage}{20cm}
\begin{bclogo}[
couleur=brown!40,
arrondi=0,
ombre=true,
couleurOmbre=gray!50,
logo=\bcplume,
noborder=true]
{\textrm{{\Large Test 3.2:
\large Statistika, kombinatorika, verjetnost
\hfill PTI-2}}}
\end{bclogo}
\end{minipage}

\vspace{0.2cm}

\begin{minipage}{20cm}
\begin{bclogo}[
couleur=brown!40,
ombre=true,
couleurOmbre=gray!50,
arrondi=0,
logo=\bcplume,
noborder=true]
{\textsc{\large Ime in priimek:
\vspace{0.2cm}
\rule{15cm}{0.2mm}}}
\end{bclogo}
\end{minipage}

\vspace{0.5cm}

%%%%%%%%%%%%%%%%%%%%%%%%%%%%%%%%%%%%%%%%%%%%%%%%%%%%%%%%%%

\def\tockeA{2+2+2}

\begin{naloga}
[\hfill $\tockeA$]

\addtocounter{tocke}{\tockeA}

V razredu je bilo na testu doseženih naslednjih 20 ocen:

\[
3,4,2,5,4,3,3,5,4,2,4,5,3,4,4,2,5,3,4,5
\]

a) Sestavi frekvenčno tabelo.

b) Izračunaj relativne frekvence.

c) Določi modus in aritmetično sredino.

\end{naloga}

\vspace{0.3cm}

%%%%%%%%%%%%%%%%%%%%%%%%%%%%%%%%%%%%%%%%%%%%%%%%%%%%%%%%%%
\newpage
\def\tockeB{2+2+2}

\begin{naloga}
[\hfill $\tockeB$]

\addtocounter{tocke}{\tockeB}

Iz števk

\[
1,\;2,\;3,\;4
\]

sestavljamo štirimestna števila brez ponavljanja števk.

a) Koliko različnih števil lahko sestavimo?

b) Koliko jih je sodih?

c) Koliko jih je večjih od 3000?

\end{naloga}

\vspace{0.3cm}

%%%%%%%%%%%%%%%%%%%%%%%%%%%%%%%%%%%%%%%%%%%%%%%%%%%%%%%%%%
\newpage
\def\tockeC{2+2+2}

\begin{naloga}
[\hfill $\tockeC$]

\addtocounter{tocke}{\tockeC}

V škatli so 4 rdeče, 3 modre in 2 zeleni kroglici.

Naključno izberemo eno kroglico.

a) Izračunaj verjetnost, da je izbrana kroglica modra.

b) Izračunaj verjetnost, da ni zelena.

c) Izračunaj verjetnost, da je rdeča ali zelena.

\end{naloga}



\vfill

\begin{kriterij*} \begin{center} \renewcommand{\arraystretch}{2.5}

\begin{tabular}{|>{\columncolor{gray!20}}c||c|c|c|c|c||c|c|c|} \hline Odstotki & $[0,45)$ & [0,45) & [0,45) & [0,45) & [0,45)& Število možnih točk& Dosežene točke & \phantom{aa}OCENA\phantom{aa} \\ \hline Ocena & 1 & 2 & 3 & 4 & 5 &\cellcolor{gray!20} $\thetocke$&\cellcolor{gray!20} \phantom{ocena} & \cellcolor{gray!20} \phantom{aaaa} \\ \hline \end{tabular} \end{center} \vspace{0.3cm}  \end{kriterij*} 

\end{document}


